% !TeX program = xelatex
\documentclass[withoutpreface,bwprint]{cumcmthesis}
\usepackage{cumcm2026}
\usepackage{xurl}
\normalem
\setCJKmainfont{SimSun}[BoldFont=SimHei,ItalicFont=KaiTi]
\setCJKfamilyfont{zhsong}{SimSun}[BoldFont=SimHei,ItalicFont=KaiTi]
\setCJKfamilyfont{zhhei}{SimHei}
\setlength{\emergencystretch}{2em}
\title{AI工具使用详情}
\hypersetup{hidelinks,pdfauthor={}}
\begin{document}
\section*{AI工具使用详情}
\subsection*{工具信息}
\noindent\begin{tabularx}{\textwidth}{@{}XX@{}}
\toprule 项目 & 填写内容 \\ \midrule
工具名称 & OpenAI Codex \\
版本或型号 & 26.908.40834 \\
使用日期 & 2026-09-10 至 2026-09-13 \\
\bottomrule
\end{tabularx}
\par\medskip
\noindent\begin{tabularx}{\textwidth}{@{}XX@{}}
\toprule 项目 & 填写内容 \\ \midrule
工具名称 & OpenAI GPT \\
版本或型号 & \texttt{gpt-6-astra} \\
使用日期 & 2026-09-10 至 2026-09-13 \\
\bottomrule
\end{tabularx}

\subsection*{具体使用目的和环节}
本参赛队完成核心建模与分析，包括模型假设、数学关系、方法选择及结果解释。AI用于以下辅助环节。
\begin{enumerate}[label=\arabic*)]
\item \textbf{数据处理。}输入题目PDF、负载与光伏数据、电价及结果模板；AI编写数据对齐、单位转换与审计程序，检查十分钟区间和模板时间标签。对应论文数据说明、问题一至四的输入及附录C。
\item \textbf{程序实现与调试。}将参赛队建立的模型、约束、参数和计算步骤作为输入，使用AI编写求解、预测与反馈控制程序，检查代码实现与数学表达的一致性，排查参数传递和边界条件实现问题。对应论文第五节及附录C。
\item \textbf{计算检查。}输入程序、结果及检查要求，使用AI编写和执行索引、费用、SOC及信息可用时点的检查程序，汇总对照与敏感性计算结果，供参赛队分析。对应论文第五、六、七节。
\item \textbf{文献检索与整理。}输入方法关键词和待支撑的论述；AI检索方法来源、整理阅读范围与引用，形成文献笔记。对应论文方法说明及参考文献；未通读文献的细节不作为已核实结论。
\item \textbf{论文与结果文件。}AI根据参赛队已经形成的分析内容、计算结果和文字草稿，辅助进行文字整理、语言修改、LaTeX排版及表图生成；模型论证、结果解释及结论由参赛队确定，并由参赛队逐项审核。对应全文、附录与结果表。
\end{enumerate}

\clearpage
\subsection*{主要提示方式与使用过程}
通过自然语言交互使用AI工具，读取题面与附件，生成、修改和执行程序，检索资料并生成与检查文档。
主要采用“给出模型与计算要求—生成代码—运行检查—定位错误—修订实现”的迭代方式。提示中限定单位、时间对齐、可用信息、SOC连续承接以及结果文件格式；发现问题后要求回到源码和原始数据核实，再更新结果及文档。

主要提示内容概括如下：
\begin{enumerate}[label=\arabic*)]
\item \textbf{代码实现：}给出模型公式、变量、约束与计算流程，要求生成相应程序，读取附件数据并按模板输出结果表。
\item \textbf{程序检查：}要求检查费用计算、可用信息时点和SOC递推的实现，输出错误位置及复现步骤；对计算中的异常提供排查建议。
\item \textbf{修订验证：}修复已确认的问题并重新验算，保留负向实验，不把可行性等同于全局最优。形成主结果、修订实验、补充实验与验证记录之间的对应关系。
\item \textbf{文字与排版：}提供参赛队的分析内容、计算结果和修改要求，使用AI整理文字、生成表图并调整LaTeX排版，编译检查公式显示、表格和源码附录。
\end{enumerate}

\subsection*{采纳、人工修改与核验}
本参赛队完成模型假设的确定、数学关系的建立、求解方法的选择和结果分析。AI输出的采纳对象主要为实现上述模型的程序、数据处理结果以及文字和排版建议。各项程序修订与计算核验情况如下。

\noindent\textbf{1. 数据处理程序。}
采纳AI生成的数据读取、时间对齐和单位转换代码，以题目规定的十分钟时段及功率、电量关系为检查依据。对结果模板中整体偏移十分钟的时间标签作了修正，并检查预测程序是否仅使用决策时点已经获得的数据。通过索引对照、物理约束检查和未来数据篡改试验，检验程序实现与信息条件的一致性。相关程序见论文附录C。

\noindent\textbf{2. 模型的程序实现。}
参赛队建立的目标函数、约束和边界条件是代码实现的依据，采纳的是AI生成并经修订的求解代码。程序修订包括将问题一日末SOC条件实现为显式等式，以及修复效率敏感性参数的传递。采用LP与MILP计算结果交叉比较；修订后总费用未变，但部分购电时段改变，因此同步更新结果表。模型及求解分析见论文第五节，计算程序见附录C。

\clearpage
\noindent\textbf{3. 对照计算与结果分析。}
采纳AI生成的对照和敏感性计算程序，用于汇总不同条件下的费用及储能运行结果。计算中对齐可用信息和初始库存，并分别实施逐笔收费下的重新计费与重新优化。参赛队负责解释这些计算结果及其适用范围：整套更新策略的收益不能全部归因于新预报；未优于主方案的对照结果予以保留；问题二至四不作全局最优性结论，同年再分析不作为未见年度验证。相关论述见论文第五、六节。

\noindent\textbf{4. 结果表与论文数值。}
采纳AI生成的结果汇总和表格排版代码。重新运行求解程序后，比较各时段购电量、充放电量、SOC和费用，并核对结果表与论文数值的精度及舍入方式。检查对象包括论文第五节的主要结果以及附录中的指定日期补充结果，确保表格展示与计算输出一致。

\noindent\textbf{5. 文献与文字整理。}
采纳与研究方法有关的资料整理及文字、排版建议。模型含义与结果解释由参赛队确定，文字调整围绕这些内容展开。引用核对以所支撑的具体论述为对象，未核查的文献细节不用于支持结论。相关引用见论文方法说明及参考文献。

\subsection*{人工核验清单}
\begin{enumerate}[label=\arabic*)]
\item 核对公式推导、单位、边界条件和符号一致性；
\item 在独立环境中运行全部代码并复现论文表图；
\item 核验引用、数据来源和事实陈述；
\item 确认核心建模与分析由参赛队主导；
\item 确认论文声明与本详情文件一致。
\end{enumerate}
\end{document}
